% Options for packages loaded elsewhere
\PassOptionsToPackage{unicode}{hyperref}
\PassOptionsToPackage{hyphens}{url}
%
\documentclass[
  10pt,
]{article}
\usepackage{amsmath,amssymb}
\usepackage{iftex}
\ifPDFTeX
  \usepackage[T1]{fontenc}
  \usepackage[utf8]{inputenc}
  \usepackage{textcomp} % provide euro and other symbols
\else % if luatex or xetex
  \usepackage{unicode-math} % this also loads fontspec
  \defaultfontfeatures{Scale=MatchLowercase}
  \defaultfontfeatures[\rmfamily]{Ligatures=TeX,Scale=1}
\fi
\usepackage{lmodern}
\ifPDFTeX\else
  % xetex/luatex font selection
\fi
% Use upquote if available, for straight quotes in verbatim environments
\IfFileExists{upquote.sty}{\usepackage{upquote}}{}
\IfFileExists{microtype.sty}{% use microtype if available
  \usepackage[]{microtype}
  \UseMicrotypeSet[protrusion]{basicmath} % disable protrusion for tt fonts
}{}
\makeatletter
\@ifundefined{KOMAClassName}{% if non-KOMA class
  \IfFileExists{parskip.sty}{%
    \usepackage{parskip}
  }{% else
    \setlength{\parindent}{0pt}
    \setlength{\parskip}{6pt plus 2pt minus 1pt}}
}{% if KOMA class
  \KOMAoptions{parskip=half}}
\makeatother
\usepackage{xcolor}
\usepackage[margin=0.80in]{geometry}
\usepackage{color}
\usepackage{fancyvrb}
\newcommand{\VerbBar}{|}
\newcommand{\VERB}{\Verb[commandchars=\\\{\}]}
\DefineVerbatimEnvironment{Highlighting}{Verbatim}{commandchars=\\\{\}}
% Add ',fontsize=\small' for more characters per line
\usepackage{framed}
\definecolor{shadecolor}{RGB}{248,248,248}
\newenvironment{Shaded}{\begin{snugshade}}{\end{snugshade}}
\newcommand{\AlertTok}[1]{\textcolor[rgb]{0.94,0.16,0.16}{#1}}
\newcommand{\AnnotationTok}[1]{\textcolor[rgb]{0.56,0.35,0.01}{\textbf{\textit{#1}}}}
\newcommand{\AttributeTok}[1]{\textcolor[rgb]{0.13,0.29,0.53}{#1}}
\newcommand{\BaseNTok}[1]{\textcolor[rgb]{0.00,0.00,0.81}{#1}}
\newcommand{\BuiltInTok}[1]{#1}
\newcommand{\CharTok}[1]{\textcolor[rgb]{0.31,0.60,0.02}{#1}}
\newcommand{\CommentTok}[1]{\textcolor[rgb]{0.56,0.35,0.01}{\textit{#1}}}
\newcommand{\CommentVarTok}[1]{\textcolor[rgb]{0.56,0.35,0.01}{\textbf{\textit{#1}}}}
\newcommand{\ConstantTok}[1]{\textcolor[rgb]{0.56,0.35,0.01}{#1}}
\newcommand{\ControlFlowTok}[1]{\textcolor[rgb]{0.13,0.29,0.53}{\textbf{#1}}}
\newcommand{\DataTypeTok}[1]{\textcolor[rgb]{0.13,0.29,0.53}{#1}}
\newcommand{\DecValTok}[1]{\textcolor[rgb]{0.00,0.00,0.81}{#1}}
\newcommand{\DocumentationTok}[1]{\textcolor[rgb]{0.56,0.35,0.01}{\textbf{\textit{#1}}}}
\newcommand{\ErrorTok}[1]{\textcolor[rgb]{0.64,0.00,0.00}{\textbf{#1}}}
\newcommand{\ExtensionTok}[1]{#1}
\newcommand{\FloatTok}[1]{\textcolor[rgb]{0.00,0.00,0.81}{#1}}
\newcommand{\FunctionTok}[1]{\textcolor[rgb]{0.13,0.29,0.53}{\textbf{#1}}}
\newcommand{\ImportTok}[1]{#1}
\newcommand{\InformationTok}[1]{\textcolor[rgb]{0.56,0.35,0.01}{\textbf{\textit{#1}}}}
\newcommand{\KeywordTok}[1]{\textcolor[rgb]{0.13,0.29,0.53}{\textbf{#1}}}
\newcommand{\NormalTok}[1]{#1}
\newcommand{\OperatorTok}[1]{\textcolor[rgb]{0.81,0.36,0.00}{\textbf{#1}}}
\newcommand{\OtherTok}[1]{\textcolor[rgb]{0.56,0.35,0.01}{#1}}
\newcommand{\PreprocessorTok}[1]{\textcolor[rgb]{0.56,0.35,0.01}{\textit{#1}}}
\newcommand{\RegionMarkerTok}[1]{#1}
\newcommand{\SpecialCharTok}[1]{\textcolor[rgb]{0.81,0.36,0.00}{\textbf{#1}}}
\newcommand{\SpecialStringTok}[1]{\textcolor[rgb]{0.31,0.60,0.02}{#1}}
\newcommand{\StringTok}[1]{\textcolor[rgb]{0.31,0.60,0.02}{#1}}
\newcommand{\VariableTok}[1]{\textcolor[rgb]{0.00,0.00,0.00}{#1}}
\newcommand{\VerbatimStringTok}[1]{\textcolor[rgb]{0.31,0.60,0.02}{#1}}
\newcommand{\WarningTok}[1]{\textcolor[rgb]{0.56,0.35,0.01}{\textbf{\textit{#1}}}}
\usepackage{graphicx}
\makeatletter
\def\maxwidth{\ifdim\Gin@nat@width>\linewidth\linewidth\else\Gin@nat@width\fi}
\def\maxheight{\ifdim\Gin@nat@height>\textheight\textheight\else\Gin@nat@height\fi}
\makeatother
% Scale images if necessary, so that they will not overflow the page
% margins by default, and it is still possible to overwrite the defaults
% using explicit options in \includegraphics[width, height, ...]{}
\setkeys{Gin}{width=\maxwidth,height=\maxheight,keepaspectratio}
% Set default figure placement to htbp
\makeatletter
\def\fps@figure{htbp}
\makeatother
\setlength{\emergencystretch}{3em} % prevent overfull lines
\providecommand{\tightlist}{%
  \setlength{\itemsep}{0pt}\setlength{\parskip}{0pt}}
\setcounter{secnumdepth}{-\maxdimen} % remove section numbering
\usepackage{enumitem}
\setlist{nosep,leftmargin=*}
\ifLuaTeX
  \usepackage{selnolig}  % disable illegal ligatures
\fi
\usepackage{bookmark}
\IfFileExists{xurl.sty}{\usepackage{xurl}}{} % add URL line breaks if available
\urlstyle{same}
\hypersetup{
  pdftitle={Author Contributions Checklist},
  hidelinks,
  pdfcreator={LaTeX via pandoc}}

\title{Author Contributions Checklist}
\usepackage{etoolbox}
\makeatletter
\providecommand{\subtitle}[1]{% add subtitle to \maketitle
  \apptocmd{\@title}{\par {\large #1 \par}}{}{}
}
\makeatother
\subtitle{Surrogate-Assisted Local Conditional Distribution Inference
with Scarce Reference Outcomes}
\author{}
\date{\vspace{-2.5em}25 August 2026}

\begin{document}
\maketitle

This form documents the data and code supporting the computational
findings in the article and describes how to reproduce those findings.

\section{Part 1: Data}\label{part-1-data}

\begin{itemize}
\item[$\square$]
  This paper does not involve analysis of external data.
\item[$\boxtimes$]
  I certify that the authors have legitimate access to and permission to
  use the data analyzed in this manuscript.
\end{itemize}

\subsection{Abstract}\label{abstract}

The empirical analyses use two public data sources and one
author-generated synthetic benchmark. The EPA study analyzes 2,814
collocated air-sensor site-days. The VitalDB study analyzes one matched
pulse-oximetry and arterial-blood-gas pair for each of 3,178 subjects.
The synthetic study contains 6,000 evaluation dossiers, of which 5,997
have valid automated scores; the primary analysis uses 2,998 cases from
Corpus A. Simulation outputs and derived data supporting all reported
figures and tables accompany the submission.

\subsection{Availability}\label{availability}

\begin{itemize}
\tightlist
\item[$\boxtimes$]
  Data \textbf{are} publicly available.
\item[$\square$]
  Data \textbf{cannot be made} publicly available.
\end{itemize}

\subsubsection{Publicly available data}\label{publicly-available-data}

\begin{itemize}
\item[$\boxtimes$]
  Data are available online at:

  \begin{itemize}
  \tightlist
  \item
    EPA, \emph{Long term performance of six air sensors across the
    United States}: \url{https://doi.org/10.23719/1531918} and
    \url{https://catalog.data.gov/dataset/dataset-long-term-performance-of-six-air-sensors-across-the-united-states}.
  \item
    VitalDB: \url{https://vitaldb.net/} and its public Web API,
    \url{https://api.vitaldb.net/}. The data resource is described in
    \url{https://doi.org/10.1038/s41597-022-01411-5}.
  \end{itemize}
\item[$\boxtimes$]
  Data are available as part of the paper's supplementary material.

  The anonymous computational supplement contains the author-generated
  synthetic dossiers, expert outcomes, parsed automated-score tables,
  derived application data, and numerical summaries used to produce the
  manuscript tables and figures.
\item[$\square$]
  Data are publicly available by request.
\item[$\square$]
  Data are available through another mechanism.
\end{itemize}

\subsubsection{Non-publicly available
data}\label{non-publicly-available-data}

Not applicable to the data required to reproduce the reported findings.
Raw API response caches and API credentials are not required for
reproduction and are not distributed. The parsed automated scores used
in the analyses are provided.

\subsection{Description}\label{description}

\subsubsection{File format(s)}\label{file-formats}

\begin{itemize}
\tightlist
\item[$\boxtimes$]
  CSV or other plain text.
\item[$\boxtimes$]
  Software-specific binary format: R data files (\texttt{.rds}).
\item[$\square$]
  Standardized binary format such as netCDF or HDF5.
\item[$\boxtimes$]
  Other: JSON, JSONL, and YAML files are used for synthetic dossiers,
  configuration, and provenance records.
\end{itemize}

\subsubsection{Data dictionary}\label{data-dictionary}

\begin{itemize}
\item[$\boxtimes$]
  Provided by the authors in the following files:

  \begin{itemize}
  \tightlist
  \item
    \texttt{data-raw/README.md};
  \item
    \texttt{data-raw/epa\_long\_term/README.md};
  \item
    \texttt{analysis/README.md};
  \item
    \texttt{reproducibility/reference\_mask\_split\_map.md};
  \item
    \texttt{reproducibility/table\_results\_map.csv};
  \item
    \texttt{reproducibility/application\_story/source\_manifest.csv}.
  \end{itemize}
\item[$\square$]
  Data files are self-describing standardized binary files.
\item[$\square$]
  A separate data dictionary is available at another URL.
\end{itemize}

\subsubsection{Additional Information}\label{additional-information}

The EPA source archive is retrieved by the supplied analysis code rather
than redistributed. VitalDB data are retrieved through the public
VitalDB interface under its access terms. The supplied synthetic
dossiers contain no real-person records. The computational supplement
contains no API credential.

\section{Part 2: Code}\label{part-2-code}

\subsection{Abstract}\label{abstract-1}

The supplied R package and analysis scripts implement the proposed
local-CDF estimators, adaptive borrowing criterion, simulation designs,
multiplier procedures, empirical analyses, and generation of manuscript
tables and figures. Python code retrieves and matches the public VitalDB
measurements. Shell and R entry points validate the numerical evidence
and rebuild all manuscript displays. Unit tests cover the estimator,
sample splitting, simulation design, application helpers, and
claim-to-result correspondence.

\subsection{Description}\label{description-1}

\subsubsection{Code format(s)}\label{code-formats}

\begin{itemize}
\tightlist
\item[$\boxtimes$]
  Script files

  \begin{itemize}
  \tightlist
  \item[$\boxtimes$]
    R
  \item[$\boxtimes$]
    Python
  \item[$\square$]
    Julia
  \item[$\square$]
    MATLAB
  \item[$\square$]
    Other
  \end{itemize}
\item[$\boxtimes$]
  Package

  \begin{itemize}
  \tightlist
  \item[$\boxtimes$]
    R
  \item[$\square$]
    Python
  \item[$\square$]
    Julia
  \item[$\square$]
    MATLAB toolbox
  \item[$\square$]
    Other
  \end{itemize}
\item[$\square$]
  Reproducible report
\item[$\boxtimes$]
  Shell script
\item[$\boxtimes$]
  Other: LaTeX manuscript and Supplement sources.
\end{itemize}

\subsubsection{Supporting software
requirements}\label{supporting-software-requirements}

\paragraph{Primary software used}\label{primary-software-used}

\begin{itemize}
\tightlist
\item
  R 4.3.2;
\item
  Python 3.12.1;
\item
  a POSIX-compatible shell;
\item
  TeX Live 2024 and Pandoc 3.1.12.3 for rebuilding the manuscript PDFs.
\end{itemize}

\paragraph{Libraries and dependencies used by the
code}\label{libraries-and-dependencies-used-by-the-code}

R package versions are recorded in \texttt{renv.lock},
\texttt{reproducibility/session\_info.txt}, and
\texttt{reproducibility/r\_package\_inventory.csv}. The package-level
requirements are listed in \texttt{DESCRIPTION}. Python packages used
for figure production are pinned in
\texttt{reproducibility/requirements-figures.txt}.

\subsubsection{Supporting system/hardware
requirements}\label{supporting-systemhardware-requirements}

The release was validated on Apple Silicon running macOS 15.4.1. No GPU
is required. The presentation-only rebuild runs on a standard desktop.
Full Monte Carlo recomputation benefits from multiple CPU cores and
substantial elapsed time but does not require a computing cluster.

\subsubsection{Parallelization used}\label{parallelization-used}

\begin{itemize}
\tightlist
\item[$\square$]
  No parallel code used.
\item[$\boxtimes$]
  Multi-core parallelization on a single machine/node.

  \begin{itemize}
  \tightlist
  \item
    Number of cores used: 4 for the manuscript-facing application and
    threshold-level recomputations; full simulation scripts accept a
    user-specified positive core count.
  \end{itemize}
\item[$\square$]
  Multi-machine/multi-node parallelization.
\end{itemize}

\subsubsection{License}\label{license}

\begin{itemize}
\tightlist
\item[$\boxtimes$]
  MIT License.
\item[$\square$]
  BSD.
\item[$\square$]
  GPL v3.0.
\item[$\square$]
  Creative Commons.
\item[$\square$]
  Other.
\end{itemize}

\subsubsection{Additional information}\label{additional-information-1}

The DeepSeek scoring calls are not required for recomputation. The valid
automated scores used by the manuscript are supplied, and the AI-stage
scripts reuse those scores without making an external API call.

\section{Part 3: Reproducibility
workflow}\label{part-3-reproducibility-workflow}

\subsection{Scope}\label{scope}

The provided workflow reproduces:

\begin{itemize}
\tightlist
\item[$\boxtimes$]
  Any numbers provided in text in the paper.
\item[$\boxtimes$]
  The computational methods presented in the paper.
\item[$\boxtimes$]
  All tables and figures in the paper.
\item[$\square$]
  Selected tables and figures only.
\end{itemize}

\subsection{Workflow}\label{workflow}

\subsubsection{Location}\label{location}

\begin{itemize}
\tightlist
\item[$\boxtimes$]
  As part of the paper's supplementary material.
\item[$\square$]
  In a Git repository disclosed during anonymous review.
\item[$\square$]
  Other.
\end{itemize}

The anonymous supplementary archive is self-contained for source
inspection, tests, reconstruction of manuscript-facing tables and
figures, and staged recomputation. A public repository location can be
added to the accepted version of this form after peer review without
changing the workflow.

\subsubsection{Format(s)}\label{formats}

\begin{itemize}
\tightlist
\item[$\boxtimes$]
  Single master code file.
\item[$\boxtimes$]
  Wrapper shell script(s).
\item[$\square$]
  Self-contained Quarto, R Markdown, or Jupyter notebook.
\item[$\boxtimes$]
  Text file that documents the workflow.
\item[$\boxtimes$]
  Makefile.
\item[$\boxtimes$]
  Other: an R package with unit tests and staged analysis scripts.
\end{itemize}

\subsubsection{Instructions}\label{instructions}

From the root of the uncompressed anonymous supplementary archive:

\begin{enumerate}
\def\labelenumi{\arabic{enumi}.}
\item
  Review \texttt{reproducibility/ANONYMOUS\_REVIEW\_README.md} and
  \texttt{renv.lock}.
\item
  Restore the R environment if desired with \texttt{renv::restore()} and
  install the Python figure dependencies from
  \texttt{reproducibility/requirements-figures.txt}.
\item
  Rebuild every main figure, the fixed-sample Supplement figure, the
  three main tables, and Supplement Tables S5--S6 with:

\begin{Shaded}
\begin{Highlighting}[]
\FunctionTok{bash}\NormalTok{ reproducibility/rebuild\_submission\_artifacts.sh}
\end{Highlighting}
\end{Shaded}
\item
  Run the implementation, design-correspondence, and claim-to-evidence
  checks with:

\begin{Shaded}
\begin{Highlighting}[]
\ExtensionTok{Rscript}\NormalTok{ reproducibility/run\_correspondence\_gate.R}
\end{Highlighting}
\end{Shaded}
\item
  Inspect the complete ordered recomputation plan without changing any
  result file with:

\begin{Shaded}
\begin{Highlighting}[]
\ExtensionTok{Rscript}\NormalTok{ reproduce.R 4 all}
\end{Highlighting}
\end{Shaded}
\item
  Full recomputation is staged by simulation, application, and AI
  components. The exact execution and finalization commands are
  documented in \texttt{docs/CANONICAL\_RECOMPUTATION\_PLAN.md}. The AI
  stage uses the supplied automated scores and does not call the
  external scoring service. After a full recomputation,
  \texttt{reproducibility/validate\_frozen\_results.R} checks the
  complete result store.
\end{enumerate}

\subsubsection{Expected run-time}\label{expected-run-time}

Approximate time needed to reproduce the analyses on a standard desktop
machine:

\begin{itemize}
\tightlist
\item[$\square$]
  \textless{} 1 minute.
\item[$\square$]
  1-10 minutes.
\item[$\square$]
  10-60 minutes.
\item[$\square$]
  1-8 hours.
\item[$\boxtimes$]
  \textgreater{} 8 hours.
\item[$\square$]
  Not feasible to run on a desktop machine.
\end{itemize}

The display rebuild in Step 3 completes in less than one minute in the
validated environment. Full Monte Carlo recomputation takes more than
eight hours and is organized into independent stages. The supplied
numerical results allow every reported table and figure to be rebuilt
without rerunning the full Monte Carlo program.

\subsubsection{Additional information}\label{additional-information-2}

\texttt{README.md}, \texttt{reproducibility/README.md},
\texttt{docs/CANONICAL\_RECOMPUTATION\_PLAN.md}, and the package tests
document the workflow and provide smaller checks of the proposed
methods. The anonymous archive omits author identities, affiliations,
private local paths, and API credentials; the complete unblinded archive
is retained separately for the editors.

\end{document}
